    \chapter{Introdução}

    \section{Contexto}
    Os Modelos de Linguagem de Grande Escala (\textit{Large Language Models -- LLMs}) \cite{ibm_llm_topics} tornaram-se componentes relevantes no desenvolvimento de aplicações baseadas em inteligência artificial. Um LLM é um modelo treinado com grandes volumes de texto, capaz de interpretar linguagem natural, gerar respostas e apoiar tarefas como sumarização, classificação, programação assistida, pesquisa semântica. Uma funcionalidade também oferecida por estes modelos é o \textit{tool calling}\cite{ibm_tool_calling}\cite{toolCallingCycle}. O \textit{tool calling} permite que os modelos não se limitem ao seu conhecimento estático, permitindo-lhes descrever um pedido de interação com sistemas externos através de uma resposta estruturada. Este pedido obedece a um formato específico, definido previamente no momento em que as ferramentas são enviadas à API de inferência.

    O ciclo de vida de um modelo divide-se em duas fases principais, treino e inferência. O treino corresponde ao processo de aprendizagem a partir de grandes conjuntos de dados, exigindo elevados recursos computacionais. A inferência corresponde à utilização do modelo já treinado para gerar respostas a novos pedidos. Atualmente, a maioria das aplicações integra LLMs através de APIs de inferência disponibilizadas por fornecedores externos. A comunicação com estas APIs é normalmente realizada usando o protocolo HTTP com representação em JSON, tanto nos pedidos como nas respostas. Antes de serem processados pelos modelos, os textos são convertidos em \textit{tokens}, pequenas unidades de informação que permitem ao modelo representar e manipular sequências de entrada e saída \cite{nvidia_tokens}. Como a resposta é gerada de forma incremental, muitas APIs suportam \textit{streaming}, permitindo enviar fragmentos da resposta ao cliente à medida que são produzidos.

    Existem vários fornecedores relevantes, incluindo \textbf{OpenAI} \cite{openai_api}, \textbf{Anthropic} \cite{anthropic_api}, \textbf{Google Gemini} \cite{google_gemini_api} e \textbf{Groq} \cite{groq_api}. Embora todos disponibilizem APIs HTTP e estruturas JSON, cada fornecedor define interfaces próprias, campos específicos, regras de autenticação e modelos distintos para \textit{tool calling}. A integração direta com vários fornecedores obriga, por isso, as aplicações a conhecerem detalhes internos de cada API.

    Ao mesmo tempo, o uso de LLMs evoluiu de interações simples para sistemas baseados em agentes (\textit{agentic workflows}) \cite{ibm_agentic_workflows}. Nestes sistemas, o modelo não se limita a devolver texto, pode também devolver instruções para chamadas a ferramentas externas, interpretar resultados intermédios e produzir uma resposta final com base nesses dados. O \textbf{Model Context Protocol (MCP)} surgiu como proposta de normalização para a comunicação entre modelos, aplicações e ferramentas externas \cite{mcp_protocol}.

    \section{Motivação e Definição do Problema}

    As Figuras \ref{fig:p1}, \ref{fig:p2} e \ref{fig:p3} ilustram o acoplamento direto entre clientes e fornecedores de inferência, bem como a dificuldade que os clientes tem em trocar de fornecedor devido a todo acoplamento que existe.

    \begin{figure}[H]
        \centering
        \begin{tikzpicture}[
            node distance=3.5cm,
            box/.style={rectangle, draw, rounded corners, align=center, text width=3.8cm, minimum height=1cm, inner sep=5pt, fill=white, thick},
            keybox/.style={rectangle, draw, rounded corners=1pt, align=center, text width=1.5cm, minimum height=0.5cm, inner sep=2pt, fill=gray!20, font=\tiny, thick},
            arrow/.style={-Latex, thick}
        ]
            % --- CLIENTE ---
            \node[box] (cliente) {App Baseada em Gen AI};
            \node[keybox, anchor=north east] (key) at (cliente.south east) {API Key};

            % --- FORNECEDOR ---
            \node[box, right=of cliente] (fornecedor) {API do \textbf{Gemini} de acesso a inferência };

            % --- LIGAÇÃO ---
            \draw[arrow] (cliente.east) -- (fornecedor.west);

        \end{tikzpicture}
        \caption{Acoplamento direto entre cliente e um fornecedor de inferência}
        \label{fig:p1}
    \end{figure}

    \begin{figure}[H]
        \centering
        \begin{tikzpicture}[
            node distance=0.8cm and 3.5cm,
            box/.style={rectangle, draw, rounded corners, align=center, text width=3.8cm, minimum height=1cm, inner sep=5pt, fill=white, thick},
            keybox/.style={rectangle, draw, rounded corners=1pt, align=center, text width=1.5cm, minimum height=0.5cm, inner sep=2pt, fill=gray!20, font=\tiny, thick},
            arrow/.style={-Latex, thick}
        ]
            % --- CLIENTE A ---
            \node[box] (clienteA) {App Baseada em Gen AI};
            \node[keybox, anchor=north east] (keyA) at (clienteA.south east) {API Key};

            % --- CLIENTE B ---
            \node[box, below=of clienteA, yshift=-0.4cm] (clienteB) {App Baseada em Gen AI};
            \node[keybox, anchor=north east] (keyB) at (clienteB.south east) {API Key };

            % --- CLIENTE C ---
            \node[box, below=of clienteB, yshift=-0.4cm] (clienteC) {App Baseada em Gen AI};
            \node[keybox, anchor=north east] (keyC) at (clienteC.south east) {API Key};

            % --- FORNECEDOR (API de acesso à inferência) ---
            % Posicionado ao lado do Cliente B para manter as setas simétricas e centradas
            \node[box, right=of clienteB] (fornecedor) {API do \textbf{Gemini} de acesso a inferência};

            % Ligações diretas e limpas das caixas dos 3 clientes para o fornecedor
            \draw[arrow] (clienteA.east) -- (fornecedor.west);
            \draw[arrow] (clienteB.east) -- (fornecedor.west);
            \draw[arrow] (clienteC.east) -- (fornecedor.west);

        \end{tikzpicture}
        \caption{Multiplicação de conexões com chaves de API descentralizadas}
        \label{fig:p2}
    \end{figure}
    \begin{figure}[H]
        \centering
        \begin{tikzpicture}[
            node distance=0.8cm and 3.5cm,
            box/.style={rectangle, draw, rounded corners, align=center, text width=3.8cm, minimum height=1cm, inner sep=5pt, fill=white, thick},
            keybox/.style={rectangle, draw, rounded corners=1pt, align=center, text width=1.5cm, minimum height=0.5cm, inner sep=2pt, fill=gray!20, font=\tiny, thick},
            cross/.style={cross out, draw=red, ultra thick, minimum size=14pt, inner sep=0pt, outer sep=0pt},
            arrow/.style={-Latex, thick}
        ]
            % --- CLIENTES (Esquerda) ---
            % Cliente A
            \node[box] (clienteA) { App Baseada em Gen AI};
            \node[keybox, anchor=north east] (keyA) at (clienteA.south east) {API Key};

            % Cliente B (Cliente 2)
            \node[box, below=of clienteA, yshift=-0.4cm] (clienteB) {App Baseada em Gen AI};
            \node[keybox, anchor=north east] (keyB) at (clienteB.south east) {API Key};

            % Cliente C
            \node[box, below=of clienteB, yshift=-0.4cm] (clienteC) {App Baseada em Gen AI};
            \node[keybox, anchor=north east] (keyC) at (clienteC.south east) {API Key};

            % --- APIS DE INFERÊNCIA (Direita) ---
            % API 1 - Alinhada na direção do Cliente A
            \node[box, right=of clienteA] (apiOpenAI) {API do \textbf{Gemini} de acesso a inferência};

            % API 2 - AGORA ALINHADA COM O CLIENTE 2 (Cliente B)
            \node[box, right=of clienteB] (apiGemini) {API da \textbf{OpenAI} de acesso a inferência};



            % Nova ligação (Cliente A aponta para a API 2 alinhada ao meio)
            \draw[arrow] (clienteA.east) -- (apiGemini.west);

        \end{tikzpicture}
        \caption{Cliente decide trocar de fornecedor}
        \label{fig:p3}
    \end{figure}

    A Figura \ref{fig:p1} apresenta um cenário em que uma aplicação comunica diretamente com um fornecedor de inferência, ficando dependente da sua API e da respetiva chave de autenticação. A Figura \ref{fig:p2} evidencia que, à medida que o número de aplicações aumenta, também aumenta o número de integrações e de chaves de API a gerir de forma independente. Finalmente, a Figura \ref{fig:p3} demonstra que a substituição do fornecedor obriga à adaptação da integração existente, consequência do forte acoplamento às interfaces específicas de cada fornecedor.

    A existência de múltiplos fornecedores de LLMs e a ausência de um protocolo comum implicam
    diferentes interfaces de acesso aos modelos, estruturas específicas de pedidos e respostas, modelos distintos de contabilização de \textit{tokens}, custos variáveis e implementações particulares para invocação de funções externas (ver Figuras \ref{fig:p1}--\ref{fig:p3}). Por exemplo, a API da \textbf{OpenAI} \cite{openai_api} organiza a interação em torno de \textit{chat completions} e \textit{tool calls}; a API da \textbf{Anthropic} \cite{anthropic_api} utiliza uma arquitetura baseada em \textit{messages} e blocos de conteúdo; e a API do \textbf{Gemini} \cite{google_gemini_api} estrutura a geração através de \textit{contents}, \textit{parts}, \textit{function declarations} e configurações de geração próprias.

    Estas diferenças criam acoplamento entre as aplicações cliente e os fornecedores, denominado de \textit{vendor lock-in}. Sempre que uma aplicação pretende trocar de modelo, suportar mais do que um fornecedor, adicionar um mecanismo de \textit{fallback} ou integrar ferramentas externas, torna-se necessário duplicar ou reimplementar lógica de tradução, autenticação, tratamento de erros e interpretação de respostas. Esta duplicação ou reimplementação aumenta a complexidade e dificulta a evolução do sistema, um problema recorrente em sistemas que combinam lógica de aplicação com integrações externas especializadas \cite{cloudzero_ai_vendor_lock_in}.

    Para além da fragmentação dos contratos, existem preocupações operacionais e de segurança. Uma aplicação que comunica diretamente com vários fornecedores precisa de gerir chaves de API, quotas, limites de utilização, métricas, auditoria, políticas de acesso e proteção contra falhas de terceiros. Sem uma camada intermédia, a implementação destas lógicas fica fragmentada por diferentes componentes, tornando mais difícil aplicar regras consistentes de autenticação, autorização, observabilidade e resiliência.

    O problema central que este projeto visa resolver é, assim, a ausência de uma camada de abstração capaz de separar as aplicações cliente das particularidades de cada fornecedor de LLMs, mantendo simultaneamente suporte para segurança, monitorização, \textit{tool calling} e integração dinâmica com ferramentas externas.

    \section{Objetivos e Proposta de Solução}
    \begin{figure}[H]
        \centering
        \begin{tikzpicture}[
            node distance=0.8cm and 2.5cm,
            box/.style={rectangle, draw, rounded corners, align=center, text width=3.5cm, minimum height=1cm, inner sep=5pt, fill=white, thick},
            gateway/.style={rectangle, draw, rounded corners, align=center, text width=3.5cm, minimum height=1.8cm, inner sep=5pt, fill=gray!20, thick},
            arrow/.style={-Latex, thick}
        ]
            % --- CLIENTES ---
            \node[box] (cliente1) {App Baseada em Gen AI};
            \node[box, below=of cliente1] (cliente2) {App Baseada em Gen AI};
            \node[box, below=of cliente2] (cliente3) {App Baseada em Gen AI};

            % --- GATEWAY ---
            \node[gateway, right=of cliente2] (gateway) {\textbf{OmniAI Gateway}};

            % --- FORNECEDORES ---
            \node[box, right=of gateway, yshift=1.2cm] (api1) {API de Inferência (\textbf{Gemini})};
            \node[box, right=of gateway, yshift=-1.2cm] (api2) {API de Inferência (\textbf{OpenAI})};

            % --- LIGAÇÕES ---
            \draw[arrow] (cliente1.east) -- (gateway.west);
            \draw[arrow] (cliente2.east) -- (gateway.west);
            \draw[arrow] (cliente3.east) -- (gateway.west);

            \draw[arrow] (gateway.east) -- (api1.west);
            \draw[arrow] (gateway.east) -- (api2.west);
        \end{tikzpicture}
        \caption{Centralização com \textbf{OmniAI Gateway}}
        \label{fig:p4_cinza}
    \end{figure}

    A solução proposta consiste numa peça de infraestrutura intermediária que se posiciona entre as aplicações baseadas em Inteligência Artificial Generativa e as diversas APIs de inferência fornecidas pelo mercado. O objetivo é centralizar o acesso aos modelos através de uma \textbf{Gateway de IA}. Para permitir que a instanciação desta \textit{Gateway} seja configurável e facilmente adaptável a diferentes cenários, a solução materializou-se na construção do \textbf{OmniAI SDK}, uma biblioteca reutilizável que abstrai a lógica de integração. A \textbf{OmniAI Gateway} instanciada a partir deste SDK atua assim como ponto de entrada para as aplicações cliente baseadas em Gen AI, centralizando o acesso a APIs de inferência, ferramentas externas e mecanismos globais de segurança e operação. Esta abordagem aproxima-se do padrão de \textbf{API Gateway} \cite{api_gateway_design} e, no contexto específico de modelos de IA, do padrão de \textbf{AI Gateway} enquanto camada centralizada de acesso a modelos \cite{ai_gateway_pattern,karl_ai_gateway}. Embora existam soluções relacionadas, como LiteLLM e OpenRouter \cite{litellm,openrouter}, o foco deste projeto está na construção de um \textbf{SDK} extensível, configurável e validado por uma \textbf{Gateway} \textit{Proof of Concept} própria.

    O \textbf{OmniAI SDK} foi desenvolvido em \textbf{Kotlin Multiplatform}, permitindo partilhar o domínio e vários componentes entre \textit{targets} diferentes e preparar a distribuição para \textbf{JVM} e \textbf{JavaScript}/\textbf{Node.js} \cite{kotlin_multiplatform}. Esta decisão torna possível disponibilizar a biblioteca em ecossistemas distintos: via \textbf{Maven} para aplicações \textbf{JVM}/\textbf{Kotlin} e via \textbf{NPM} para ambientes \textbf{Node.js}/\textbf{TypeScript}.

    A arquitetura do \textbf{OmniAI SDK} foi organizada segundo os princípios da Arquitetura Hexagonal, também conhecida como \textit{Ports and Adapters} \cite{hexagonal_ports}. O domínio comum fica isolado de contratos externos, enquanto os \textit{inbound adapters} traduzem pedidos de clientes compatíveis com \textbf{OpenAI}, \textbf{Anthropic} ou \textbf{Gemini} para o domínio comum, e os \textit{outbound adapters} convertem esse pedido já em dominio comum para o formato esperado pela API de inferência configurada, assegurando a comunicação com o fornecedor correspondente.

    Para tornar a solução configurável e escalável, foi criada uma \textit{pipeline} interna de \textit{interceptors}, inspirada no modelo de filtros e cadeia de filtros dos \textbf{Servlets} \cite{java_servlet_filter}\cite{java_servlet_filterchain}. Esta \textit{pipeline} permite injetar comportamentos específicos tanto antes como após a chamada ao fornecedor de inferência ou ao próximo \textit{interceptor} na \textit{pipeline}, incluindo autenticação, autorização, métricas, \textit{rate limiting}, \textit{circuit breaker}, \textit{fallback} e \textit{logging}.

    Para garantir que a \textbf{OmniAI Gateway} atue como uma peça de infraestrutura completa de inteligência artificial, a arquitetura incorpora um \textbf{MCP Broker}. Este módulo permite a exposição  e centralização de múltiplas ferramentas externas através do \textbf{Model Context Protocol (MCP)} \cite{mcp_protocol}.   

    \section{Estrutura do Documento}
    Este documento encontra-se organizado de forma progressiva, partindo do problema de integração e avançando até à validação da solução construída. O Capítulo 2 apresenta o estudo prévio: primeiro a análise comparativa das APIs de inferência e depois o estudo do \textbf{Model Context Protocol}. O Capítulo 3 descreve a arquitetura e desenho do sistema para suporte à inferência, incluindo a visão de alto nível e a modelação do domínio comum. O Capítulo 4 foca-se na implementação técnica e evolução do sistema, detalhando a Arquitetura Hexagonal, os vários adaptadores e portas, a \textit{pipeline} de \textit{interceptors}, a DSL, bem como os mecanismos de segurança, observabilidade e resiliência. O Capítulo 5 é dedicado exclusivamente ao \textbf{MCP Broker} e à integração com ferramentas externas, detalhando a arquitetura, desenho e implementação deste componente. O Capítulo 6 apresenta as provas de conceito efetivamente realizadas na \textbf{OmniAI Gateway} (\textit{Proof of Concept}). O Capítulo 7 sintetiza as conclusões, limitações e linhas de evolução futura. O Apêndice \ref{chap:historico_planeamento} documenta o histórico de planeamento e identifica os principais riscos considerados durante a fase de desenvolvimento.